% ============================================================================
% SUGGESTED text for /Users/talyaron/Documents/ChatGPT/מאמרים/output/pdf/freedi-arxiv/source/main.tex
% Prepared 2026-09-11 from scientific-research/2026-09-11-clustering-scale-probe.
% NOT applied. The pilot's accuracy numbers rest on PROVISIONAL (model-authored)
% labels; do not merge the accuracy sentences until the human labels are in and
% analyze.mjs has been re-run with --labels=annotation/labels.human.jsonl.
% Line numbers refer to main.tex as read on 2026-09-11 (416 lines).
% ============================================================================

% ---------------------------------------------------------------------------
% (1) After L336 (the (b_in, b_out) = (2000, 200) assumption paragraph) — ADD:
% ---------------------------------------------------------------------------
A price correction applies to the full-list curve. Re-reading the archived usage
records shows that 62.9\% of its prompt tokens were billed as prompt-cache
\emph{writes}, which the provider's 11~September~2026 schedule prices at
$1.25\times$ the input rate, and that none were cache reads. With that surcharge
the full-list projection rises from \$0.44 to \$0.50 at 500 proposals, from
\$1.62 to \$1.86 at 1,000, and from \$149 to \$172 at 10,000. The local-work
budget remains an assumption; the component measurement in
Section~\ref{sec:scaleprobe} does not replace it.

% ---------------------------------------------------------------------------
% (2) L345--346 ("No cached input tokens were reported.") — REPLACE that sentence with:
% ---------------------------------------------------------------------------
No cached-input reads were reported; the cache \emph{writes} that were reported
are billed at a premium and are included in the corrected figures above.

% ---------------------------------------------------------------------------
% (3) L348 (the 200-proposal paragraph) — REPLACE with:
% ---------------------------------------------------------------------------
Token accounting cannot establish an accuracy curve or a failure threshold at
200 proposals. Long-context findings motivate controlled tests
\citep{liu2024lost} but supply no proposal-count threshold for this task. A
first controlled probe (Section~\ref{sec:scaleprobe}) found no decline in
full-list judging accuracy between 100 and 500 proposals.

% ---------------------------------------------------------------------------
% (4) New subsection, suggested after the cost-projection subsection (before L350):
% ---------------------------------------------------------------------------
\subsection{A Budgeted Probe of Retrieve-then-Judge}\label{sec:scaleprobe}

We tested one component of the mechanism in isolation: does showing the judge
only the 15 nearest proposals by embedding cosine (B), instead of the full pool
(A), reduce cost without an evident loss in matching decisions?

\paragraph{Design.}
\begin{itemize}
  \item \textbf{Data.} Twenty held-out proposals from one public deliberation
    (Pol.is, Bowling Green, 607 approved comments, CC~BY~4.0): ten with at least
    one allowable match in the bank and ten with none.
  \item \textbf{Pools.} Each query was judged against nested pools of 100, 200,
    and 500 existing proposals, in two presentation orders.
  \item \textbf{Controls.} The same model (\texttt{gpt-5.6-luna}), rubric,
    schema, output cap, and candidate texts were used in both arms. B kept A's
    relative order.
  \item \textbf{Scale.} There were 240 paired judgements, at a total API cost of
    US\$0.27.
\end{itemize}

\paragraph{Cost.} B's input stayed at about 655 tokens while A's grew linearly
(2,734, 5,150, and 12,520 tokens). The median paired cost ratio was
$2.6\times$, $4.3\times$, and $9.4\times$ at 100, 200, and 500 proposals
(undiscounted list prices). Output tokens, similar in both arms, bound the saving.

\paragraph{Decisions.} [PROVISIONAL --- re-derive from human labels before use.]
\begin{itemize}
  \item Full-list correctness did not decline from 100 to 500 proposals (28, 31,
    and 31 of 40 cells). There were no false joins at 200 or 500.
  \item Retrieval placed an approved target among its 15 candidates for every
    match-present query at every size.
  \item The judge nevertheless declined these targets more often in B: at 500
    proposals, 6 paired cells were correct only in A and none only in B.
  \item Its stated reasons indicate a contrast effect. When every candidate is a
    near-variant, the judge demands closer equivalence.
\end{itemize}
Whether this is a loss of recall or a gain in precision depends on the adjudicated
labels. The pairwise attachment judge used in production was not tested.

\paragraph{Scope.} The probe measures a simplified component at sampled
states. It is not a streamed run of the seven-stage pipeline and not a
cluster-level F1 result. It uses one English discussion and one model; the
bootstrap intervals, resampled by query, do not cover variation across discussions.

% ---------------------------------------------------------------------------
% (5) L367 (Future evaluation, "Twenty independently labelled query proposals ...
%     requires 240 decision calls.") — REPLACE with:
% ---------------------------------------------------------------------------
The pilot described in Section~\ref{sec:scaleprobe} has been run; its labels
await independent human adjudication. The next step is a streamed comparison at
100, 200, and 500 arrivals, with at least two orders and more than one
labelled discussion. Every provider call of the complete pipeline, including
background maintenance, is to be metered under a shared budget.
